\hypertarget{simd-vectorization-with-python}{%
\chapter{SIMD Vectorization with
Python}\label{simd-vectorization-with-python}}

SIMD (Single Instruction, Multiple Data) allows a single CPU instruction
to perform the same operation on multiple values simultaneously (e.g.,
adding 4 floats in one cycle).

\hypertarget{avx-and-sse-in-the-standard-library}{%
\subsection{71.1 AVX and SSE in the Standard
Library}\label{avx-and-sse-in-the-standard-library}}

While the CPython interpreter loop doesn't use SIMD, many of its
underlying C-extensions do. * \textbf{\passthrough{\lstinline!base64!}}:
Uses SIMD to accelerate bit-shifting operations. *
\textbf{\passthrough{\lstinline!hashlib!}}: Modern SHA implementations
use hardware-accelerated instructions available on Intel (SHA-NI) and
ARM (NEON).

\hypertarget{vectorizing-with-numpy}{%
\subsection{71.2 Vectorizing with NumPy}\label{vectorizing-with-numpy}}

NumPy's universal functions (\passthrough{\lstinline!ufuncs!}) are
compiled with SIMD support. When you run
\passthrough{\lstinline!arr1 + arr2!}, the underlying C code uses vector
registers to process chunks of the arrays at once, achieving throughput
that pure Python loops can never match.

\begin{center}\rule{0.5\linewidth}{0.5pt}\end{center}
